\documentclass[11pt]{article}
\usepackage[T1]{fontenc}
\usepackage[utf8]{inputenc}
\usepackage{lmodern}
\usepackage[a4paper,margin=1in]{geometry}
\usepackage{microtype}
\usepackage{amsmath,amssymb,amsthm,mathtools}
\usepackage{booktabs,array}
\usepackage[dvipsnames]{xcolor}
\usepackage[colorlinks=true,linkcolor=MidnightBlue,citecolor=MidnightBlue,urlcolor=MidnightBlue]{hyperref}
\usepackage{enumitem}
\usepackage{setspace}
\usepackage{fancyhdr}
\usepackage{titlesec}
\setlength{\headheight}{14pt}
\pagestyle{fancy}
\fancyhf{}
\fancyhead[L]{Repaired split-forest decidability proof, v2 (M6$'$)}
\fancyhead[R]{May 2026}
\fancyfoot[C]{\thepage}
\titleformat{\section}{\large\bfseries\color{MidnightBlue}}{\thesection}{0.6em}{}
\titleformat{\subsection}{\normalsize\bfseries\color{MidnightBlue}}{\thesubsection}{0.5em}{}
\setlist[itemize]{topsep=4pt,itemsep=2pt,leftmargin=1.5em}
\setlist[enumerate]{topsep=4pt,itemsep=2pt,leftmargin=1.8em}
\onehalfspacing

\newtheorem{theorem}{Theorem}[section]
\newtheorem{proposition}[theorem]{Proposition}
\newtheorem{lemma}[theorem]{Lemma}
\newtheorem{corollary}[theorem]{Corollary}
\theoremstyle{definition}
\newtheorem{definition}[theorem]{Definition}
\newtheorem{construction}[theorem]{Construction}
\theoremstyle{remark}
\newtheorem{remark}[theorem]{Remark}
\newtheorem{example}[theorem]{Example}

\newcommand{\DR}{\mathsf{DR}}
\newcommand{\PO}{\mathsf{PO}}
\newcommand{\PP}{\mathsf{PP}}
\newcommand{\PPI}{\mathsf{PPI}}
\newcommand{\EQ}{\mathsf{EQ}}
\newcommand{\Base}{\mathsf{B}}
\newcommand{\Inv}{\operatorname{Inv}}
\newcommand{\comp}{\operatorname{comp}}
\newcommand{\ALCIRCC}[1]{\mathcal{ALCI}_{\mathsf{RCC#1}}}
\newcommand{\Vrt}{\mathsf{Vrt}}
\newcommand{\Sib}{\mathsf{Sib}}
\newcommand{\Front}{\mathsf{Front}}
\newcommand{\Mos}{\mathsf{Mos}}
\newcommand{\Ctx}{\mathsf{Ctx}}
\newcommand{\Pos}{\mathsf{Pos}}
\newcommand{\Trip}{\mathsf{Trip}}
\newcommand{\Pair}{\mathsf{Pair}}
\newcommand{\Port}{\mathsf{Port}}
\newcommand{\Mate}{\mathsf{Mate}}
\newcommand{\Reach}{\mathsf{Reach}}
\newcommand{\Side}{\mathsf{Side}}
\newcommand{\Q}{\mathcal{Q}}
\newcommand{\U}{\mathcal{U}}
\newcommand{\T}{\mathcal{T}}
\newcommand{\md}{\operatorname{md}}
\newcommand{\Cl}{\operatorname{cl}}
\newcommand{\Typ}{\operatorname{Typ}}
\newcommand{\Stratum}{\mathsf{Strat}}
\newcommand{\Subsrc}{\mathsf{subsrc}}
\newcommand{\eqc}{\equiv_{\EQ}}
\newcommand{\eqQEQ}{\equiv^{\Q}_{\EQ}}

\begin{document}

\begin{center}
{\LARGE\bfseries Repaired Split-Forest Decidability Proof for $\ALCIRCC{5}$ (v2):\\[2pt]
Recursive Verticalization (M6$'$) to Close the Round-5 Side-Witness Gap\par}
\vspace{0.5em}
{\large A delta to the 31-page round-4 repaired proof~\cite{GPT55Repaired}\par}
\vspace{0.75em}
{\normalsize Claude (Anthropic, Opus~4.7), prompted by Michael Wessel}\par
\vspace{0.45em}
{\small May~2026}
\end{center}

\begin{center}
\fbox{\parbox{0.92\textwidth}{%
\textbf{Disclaimer.}\;
This document was drafted by Claude (Anthropic), prompted by Michael
Wessel, in response to GPT-5.5's fifth formal review of the
repaired split-forest proof.  The results have not been independently
verified by human domain experts.  We invite scrutiny and
corrections.}}
\end{center}

\begin{abstract}
GPT-5.5's fifth formal review~\cite{GPT55Round5Review} of the 31-page
repaired split-forest decidability proof for $\ALCIRCC{5}$ identifies
one load-bearing mathematical defect: the verticalization discipline
(M6) of Definition~5.1 in~\cite{GPT55Repaired} is too strong.  It
forbids non-vertical $\PP/\PPI$ relations between sibling-frontier
mosaic positions, but such relations occur in satisfiable
witness-generated models, as witnessed by the stress concept
\[
  C_{\rm side}\;:=\;\exists\DR.(A\sqcap\exists\PP.B)\;\sqcap\;\exists\DR.B,
\]
which is satisfiable in abstract RCC5 (and is reported \emph{SAT} by
all three Python reasoners in the project repository), yet rejected
by the round-4 (M6).  This delta paper repairs the defect by
replacing the flat verticalization $\Vrt\subseteq P$ with a
\emph{recursive vertical-stratum forest}, in which every selected
$\DR/\PO$ side-witness may itself open a local vertical stratum
containing its own selected $\PP/\PPI$ ancestors and equality mates.
The new discipline (M6$'$) is the natural recursive lift of the
round-4 (M6); the rest of the round-4 architecture (occurrence-level
equality, port-indexed mate clusters $\Mate_\Q(C,s,p)$, the abstract
triple graph $\T_\Q$, request-closed cycles in place of
$\omega$-acceptance) carries over without further change.  The
explicit cubic side-width bound $m(n)=O(n^3)$ of the round-4 paper is
replaced by a finite recurrence over modal depth; the certificate
space is still bounded by an effectively computable
$B(C_0)\leq 2^{2^{p(n)}}$ for some computable $p$.
\end{abstract}

\tableofcontents
\bigskip

%% ==================================================================
\section{Scope and relationship to the round-4 paper}
\label{sec:scope}
%% ==================================================================

This is a delta to the 31-page round-4 paper
\cite{GPT55Repaired}.  Specifically:

\begin{itemize}
\item Section~5 (Finite interfaces and mosaics) of \cite{GPT55Repaired}
  is amended in \S\ref{sec:m6prime} below: Definition~5.1 (Mosaic) is
  restated with a recursive vertical-stratum forest, and the four
  sub-clauses of (M6) are replaced by their recursive analogues
  (M6$'$).
\item Section~7.2 (Verticalization soundness) of \cite{GPT55Repaired}
  is amended in \S\ref{sec:lifting}: Lemma~7.1 (Vertical realisation
  of $\PP/\PPI$ pairs) is restated as Lemma~\ref{lem:vertreal-recursive}
  with the recursive case analysis.
\item Section~6.5 (Validity conditions) of \cite{GPT55Repaired} is
  amended in \S\ref{sec:validity}: clauses (V3), (V4), (V5), and
  (V9.Eb) are restated to traverse the stratum forest, in addition to
  the port-indexed mate clusters introduced in round 4.
\item Section~4.3 (Explicit $\DR/\PO$ witness placement and side
  interfaces) of \cite{GPT55Repaired} is amended in
  \S\ref{sec:side-bound}: the polynomial bound $m(n)=O(n^3)$ is
  replaced by a finite recurrence over modal depth.  The certificate
  space is still finite and effectively computable.
\item All other sections of \cite{GPT55Repaired} are preserved
  verbatim.  In particular: witness thinning (\S3 of
  \cite{GPT55Repaired}), occurrence-level equality (\S4.1), the
  port-indexed mate clusters $\Mate_\Q(C,s,p)$ (\S6.3), the abstract
  triple graph $\T_\Q$ (\S5.2), Renz--Nebel-based mosaic patchwork
  (\S5.4--\S5.7), request-fulfilled cyclic repetition (\S6.5), and the
  no-automata regularization (\S8) all carry over without
  modification.
\end{itemize}

The repair does not change the soundness chain (\S7 of
\cite{GPT55Repaired} except for Lemma~7.1) and does not change the
completeness chain (\S8 of \cite{GPT55Repaired} except for the
side-interface construction at \S4.3).  In particular, the proof
remains no-automata, in the sense that the only finite-state machinery
used is the request-closed blocking cycle of Section~6.5 in
\cite{GPT55Repaired}.

%% ==================================================================
\section{The defect, with the round-5 stress concept}
\label{sec:defect}
%% ==================================================================

\subsection{The round-4 verticalization discipline (M6)}

Definition~5.1 of \cite{GPT55Repaired} declares, for every mosaic
$M=(P,\tau,L,\Vrt)$ over a finite position set $P$, that
$\Vrt\subseteq P$ is the set of \emph{vertical} positions (the
source occurrence of the mosaic, its explicit equality mates, and any
selected $\PP/\PPI$ containment witnesses recorded in the mosaic) and
$\Sib:=P\setminus\Vrt$ is the set of \emph{sibling-frontier}
positions (selected $\DR/\PO$ witnesses and ordinary sibling /
boundary nodes).  Clause (M6) then states:

\begin{enumerate}[label=(M6.\alph*)]
\item $\EQ$-closure of strata: $L(p,q)=\EQ$ implies $p\in\Vrt$ iff $q\in\Vrt$;
\item containment labels are vertical: $L(p,q)\in\{\PP,\PPI\}$ implies $p,q\in\Vrt$;
\item cross-stratum labels are $\DR/\PO$: if exactly one of $p,q$ lies in $\Vrt$, then $L(p,q)\in\{\DR,\PO\}$;
\item within sibling-frontier: if $p,q\in\Sib$ are in different equality classes and $p\neq q$, then $L(p,q)\in\{\DR,\PO\}$.
\end{enumerate}

\subsection{The stress concept}
\label{ssec:stress}

Let $A,B$ be atomic concepts of $\Cl(C_0)$ and consider
\[
  C_{\rm side}\;:=\;\exists\DR.(A\sqcap\exists\PP.B)\;\sqcap\;\exists\DR.B.
\]

\begin{example}[Concrete satisfying model]\label{ex:stress-model}
Let $\Delta=\{x,y,z\}$ with $A^{\mathcal{I}}=\{y\}$,
$B^{\mathcal{I}}=\{z\}$, and atomic relation labels
\[
  L(x,y)=\DR,\quad L(x,z)=\DR,\quad L(y,z)=\PP
  \;\;(\,L\text{ extended to inverses and reflexive }\EQ\text{ in the usual way}).
\]
JEPD and the RCC5 composition table are satisfied: the only nontrivial
triple is $(x,y,z)$ with $\comp(\DR,\PP)=\{\DR,\PO,\PP\}$, and the
chosen label $L(x,z)=\DR$ lies in $\comp(\DR,\PP)$.  The model
satisfies $C_{\rm side}$ at $x$: the first conjunct is witnessed by
$y$ (which satisfies $A\sqcap\exists\PP.B$, since $y$ is in $A$ and
$L(y,z)=\PP$ with $z\in B$), and the second conjunct is witnessed by
$z\in B$.
\end{example}

The cover-tree tableau, the baseline quasimodel reasoner, and the
cycle-aware quasimodel reasoner of the repository all report
$C_{\rm side}$ \emph{SAT}; the verification script
\texttt{verification/python/wp7\_side\_witness\_stress.py} records the
agreement.  Hence the concept is genuinely satisfiable in abstract
RCC5.

\subsection{Why round-4 (M6) rejects $C_{\rm side}$}

Consider the candidate certificate that the round-4 proof would
construct for a model of $C_{\rm side}$ rooted at $x$.  At the source
$x$ we place two selected $\DR$ witnesses, one for each conjunct, say
$p_y$ for $y\models A\sqcap\exists\PP.B$ and $p_z$ for $z\models B$.
The mosaic at $x$ has $\Vrt_M=\{x\}$ (since $x$ has no selected
$\PP/\PPI$ witnesses; the upward containment from $x$ is empty) and
$\Sib_M=\{p_y,p_z\}\cup\{\text{boundary nodes}\}$.  By
Example~\ref{ex:stress-model}, $L(p_y,p_z)=\PP$.  But clause (M6.d)
of round-4 requires $L(p_y,p_z)\in\{\DR,\PO\}$, contradiction.

So the round-4 (M6) rejects every certificate built from the
witness-generated submodel of any satisfying model of $C_{\rm side}$.
Concretely, the (M6) clause that fails is the within-sibling-frontier
clause (M6.d): it asserts that two distinct equality classes inside
$\Sib$ admit only $\{\DR,\PO\}$ labels, whereas $\PP/\PPI$ between
two $\DR$-witnesses of the source can be forced by the
witness-generated containment structure of $C_{\rm side}$.

\begin{remark}\label{rem:why-round4-overshoots}
The round-4 proof needed (M6.d) to make Lemma~7.1 (vertical
realisation of $\PP/\PPI$ pairs) of \cite{GPT55Repaired} go through:
that proof argued, in Subcase~(ii.e), that every $\PP/\PPI$ mosaic
label is between vertical positions because (M6.b) forbids
$\PP/\PPI$ across the $\Vrt$/$\Sib$ boundary and (M6.d) forbids
$\PP/\PPI$ within $\Sib$.  The recursive repair below preserves the
\emph{conclusion} of Lemma~7.1 (vertical realisation modulo
$\EQ$-mate substitution) while admitting nested local vertical
strata, so the case analysis no longer needs to exclude
sibling-frontier $\PP/\PPI$ pairs.
\end{remark}

%% ==================================================================
\section{Recursive verticalization}
\label{sec:m6prime}
%% ==================================================================

The repair is to replace the single vertical set $\Vrt\subseteq P$ by
a finite \emph{forest} of nested vertical strata.  Each stratum has
its own source, its own equality mates, and its own selected
$\PP/\PPI$ witnesses, recursively.  The forest carries the explicit
nesting and the labels between positions in distinct strata are
still constrained as in round-4 (cross-stratum $\DR/\PO$).

\subsection{Stratum forest}

\begin{definition}[Vertical stratum]\label{def:stratum}
A \emph{vertical stratum} on a finite position set $P$ is a tuple
$V=(P_V,s_V,\mu_V)$ where $P_V\subseteq P$ is the stratum's support,
$s_V\in P_V$ is its \emph{source position}, and
$\mu_V\subseteq P_V$ is the stratum's set of \emph{explicit equality
mates of $s_V$}, with $s_V\in\mu_V$.  The positions in
$P_V\setminus\mu_V$ are the stratum's \emph{selected
$\PP/\PPI$-witnesses}.
\end{definition}

\begin{definition}[Stratum forest]\label{def:stratum-forest}
A \emph{stratum forest} on a finite position set $P$ is a pair
$\mathcal V=(\Stratum,\Subsrc)$ where
\begin{enumerate}[label=(SF\arabic*)]
\item $\Stratum=\{V_0,V_1,\dots,V_k\}$ is a finite set of vertical
  strata on $P$ (Definition~\ref{def:stratum});
\item $V_0$ is distinguished as the \emph{root stratum}; its source
  $s_{V_0}$ is called the \emph{mosaic source};
\item $\Subsrc:\Stratum\setminus\{V_0\}\to P$ is the \emph{parent
  position map}: $\Subsrc(V_i)$ is the position in some other stratum
  $V_{j}$ at which $V_i$ is rooted as a nested local stratum;
\item the parent relation $\Subsrc$ induces a forest on $\Stratum$
  with $V_0$ as the global root: for every non-root stratum $V_i$,
  $\Subsrc(V_i)\in P_{V_j}$ for some uniquely determined $V_j\neq
  V_i$, and the relation $V_i\prec V_j$ iff
  $\Subsrc(V_i)\in P_{V_j}$ is acyclic;
\item supports are stratified: distinct strata have distinct sources,
  and for any two strata $V_i,V_j$, $P_{V_i}\cap P_{V_j}$ is either
  empty or contains exactly the parent position $\Subsrc(V_i)$ when
  $V_i$ is nested at a position of $V_j$.
\end{enumerate}
\end{definition}

\begin{remark}
(SF5) is the key sharing rule.  A position $p\in P$ can be a
\emph{sibling-frontier} member of $V_j$ and \emph{simultaneously} the
source of a nested stratum $V_i$ rooted at $p$.  In that case
$p\in P_{V_i}\cap P_{V_j}$ but the rest of $P_{V_i}$ does not
intersect $P_{V_j}$.
\end{remark}

We extend the stratum-membership relation pointwise:

\begin{definition}[Stratum membership]\label{def:stratum-mem}
For $p\in P$, write
\[
  \Stratum(p)\;=\;\{\,V\in\Stratum\;:\;p\in P_V\,\}.
\]
By (SF5), $|\Stratum(p)|\in\{1,2\}$, with $|\Stratum(p)|=2$ exactly
when $p$ is the parent position of some nested stratum.

We say $p$ is \emph{vertical in $V$} if $p\in P_V$.  We say a pair
$(p,q)$ is \emph{co-vertical} if there exists $V\in\Stratum$ with
$\{p,q\}\subseteq P_V$; otherwise it is \emph{cross-stratum}.  The
\emph{residual sibling frontier} of $V_j$ at a position $p\in P_{V_j}$
is the set of positions in $P_{V_j}$ that are in $V_j$ but in no
nested stratum rooted at $p$.
\end{definition}

\subsection{Restated mosaic definition}

\begin{definition}[Mosaic, v2]\label{def:mosaic-v2}
A mosaic over a finite position set $P$ is a tuple
\[
  M\;=\;(P,\tau,L,\mathcal V)
\]
where $\tau:P\to\Typ(C_0)$, $L:P^2\to\Base$, and $\mathcal V$ is a
stratum forest on $P$ (Definition~\ref{def:stratum-forest}).  The
tuple satisfies axioms (M1)--(M5) of Definition~5.1 in
\cite{GPT55Repaired}, unchanged in form, together with the recursive
verticalization discipline:
\begin{enumerate}[label=(M6$'$.\alph*)]
\item \textbf{$\EQ$-closure within strata.}  If $L(p,q)=\EQ$ and
  $V\in\Stratum$ with $p\in P_V$, then $q\in P_V$.  In other words,
  every $\EQ$-class lies entirely inside some single stratum.

\item \textbf{Containment labels are co-vertical.}  If $L(p,q)\in
  \{\PP,\PPI\}$, then $(p,q)$ is co-vertical
  (Definition~\ref{def:stratum-mem}): there exists $V\in\Stratum$ with
  $\{p,q\}\subseteq P_V$, and $L(p,q)$ matches the local containment
  data of $V$ (i.e., if $L(p,q)=\PP$ then $q$ is the source $s_V$ or
  an equality mate, and $p$ is a selected $\PPI$-witness of $V$; the
  $\PPI$ case is symmetric).

\item \textbf{Cross-stratum labels are $\DR/\PO$ or $\EQ$.}  If $(p,q)$
  is cross-stratum and $p\neq q$, then $L(p,q)\in\{\DR,\PO\}$ unless
  $(p,q)$ is identified by an explicit $\EQ$-class, in which case
  $L(p,q)=\EQ$ and (by (M6$'$.a)) both $p$ and $q$ lie inside one
  stratum.

\item \textbf{Residual sibling frontier.}  If $p,q$ are both in the
  residual sibling frontier of some common stratum $V_j$ at a common
  parent position (i.e., neither $p$ nor $q$ opens a nested stratum
  containing the other), then $L(p,q)\in\{\DR,\PO\}$ or $L(p,q)=\EQ$
  with $p,q$ in the same equality class.
\end{enumerate}
\end{definition}

The intended reading is: each $\PP/\PPI$ label in the mosaic is
\emph{justified} by some stratum in $\mathcal V$ — either the root
stratum $V_0$ (the round-4 case) or a nested stratum opened at a
sibling-frontier position of $V_0$.  Pairs of positions whose vertical
contexts are mutually incomparable still receive only $\DR/\PO$
labels, which is the local arc-consistent regime that the
patchwork-based completeness argument requires.

\begin{lemma}[Round-4 (M6) is the depth-1 case]\label{lem:round4-recovery}
If $\mathcal V=(\{V_0\},\emptyset)$ (the stratum forest has only the
root stratum and no nested strata), then Definition~\ref{def:mosaic-v2}
specializes to Definition~5.1 of \cite{GPT55Repaired} with
$\Vrt=P_{V_0}$.  In particular, every round-4 mosaic is a v2 mosaic.
\end{lemma}

\begin{proof}
Direct unfolding of definitions.  (M6$'$.a) becomes (M6.a) since the
only stratum is $V_0$ and $\Vrt=P_{V_0}$.  (M6$'$.b) becomes (M6.b)
because co-vertical reduces to ``both in $V_0$''.  (M6$'$.c) and
(M6$'$.d) become (M6.c) and (M6.d) by the same reduction.
\end{proof}

\begin{remark}\label{rem:why-recursive}
Recursive verticalization is not an arbitrary generalization: it is
exactly what the witness-generated containment structure produces.
When the witness selector chooses, at a source $u$, a $\DR$-witness
$w$ that itself has a $\PP$-witness $w'$ inside the model, then the
selector at $w$ produces a new local mosaic context whose source is
$w$.  In round-4, that nested context was \emph{separate} from $u$'s
mosaic.  The defect arose because $w'$ can also be a selected
$\DR/\PO$ witness of $u$ (when $w'$ is $\DR$ to $u$ but $\PP$ to $w$),
in which case the round-4 mosaic at $u$ records both $w$ and $w'$ at
the frontier with no place to put the forced $\PP$ label between
them.  The repair pulls $w$ and $w'$ into a shared nested stratum
$V_w$ inside $u$'s mosaic, recording the $\PP$ label as a local
co-vertical label of $V_w$.
\end{remark}

%% ==================================================================
\section{Verticalization soundness (restated Lemma~7.1)}
\label{sec:lifting}
%% ==================================================================

We lift (M6$'$) from a single mosaic to the entire unfolding $\U(\Q)$.
The argument generalizes the round-4 case analysis of
Lemma~7.1 in~\cite{GPT55Repaired} (Subcases (ii.a)--(ii.e) of the
\emph{mosaic-label} case) by adding the recursive stratum-membership
data.

\begin{lemma}[Vertical realisation, recursive]\label{lem:vertreal-recursive}
Let $\Q$ be a valid certificate satisfying (V1)--(V10) and (M6$'$) on
every mosaic in $\Mos_\Q$.  Let $u,v\in\Omega(\U(\Q))$ and
$R\in\{\PP,\PPI\}$.  Suppose $L_{\U(\Q)}(u,v)=R$ in the unfolding.
Then either
\begin{enumerate}[label=(\alph*)]
\item there is a strict vertical $\Reach_R$-chain from $u$ to $v$ in
  the generated backbone of $\U(\Q)$ (\emph{not necessarily through
  the mosaic source}; the chain may live in a nested stratum), or
\item there exist $u',v'\in\Omega(\U(\Q))$ with $u\eqc u'$ and
  $v\eqc v'$ such that $(u',v')$ falls under~(a).
\end{enumerate}
\end{lemma}

\begin{proof}
As in the proof of Lemma~7.1 of \cite{GPT55Repaired}, the label
$L_{\U(\Q)}(u,v)\in\Base$ is assigned by either (i) a vertical
generated cover edge recorded in the certificate's containment data,
or (ii) a mosaic label $L_M(p,q)$ where $u$ realises position $p$ and
$v$ realises position $q$ of some $M\in\Mos_\Q$.

\emph{Case (i): vertical generated cover.}  Unchanged from the
round-4 proof: $u$ and $v$ lie on the same vertical backbone, with
$u\Reach_R v$, by (V4)/(V5).  This is case~(a).

\emph{Case (ii): mosaic label.}  $L_M(p,q)=R\in\{\PP,\PPI\}$.  By
(M6$'$.b), there exists $V\in\Stratum_M$ with $\{p,q\}\subseteq
P_V$.  Within stratum $V$, positions are partitioned into the
source $s_V$, its explicit equality mates $\mu_V$, and selected
containment witnesses of $V$.  We argue by the same case-by-case
analysis as round-4, but \emph{within} stratum $V$:

\begin{itemize}
\item Subcase (ii.a$'$): $p=s_V$ and $q$ is a selected $R$-witness of
  $V$.  The realising pair $(u,v)$ is the source of $V$ and its
  selected $R$-witness on $V$'s vertical backbone, so $u\Reach_R v$.
  This is case~(a).

\item Subcase (ii.b$'$): $p\in\mu_V$ (an equality mate of $s_V$) and
  $q$ is a selected $R$-witness of $V$.  By (V9.Eb) applied to
  stratum $V$ (see~\S\ref{sec:validity} below), the source $s_V$ is
  realised by an occurrence $u'$ with $u\eqc u'$ and $u'\Reach_R v$.
  This is case~(b) with $v'=v$.

\item Subcase (ii.c$'$): $p=s_V$ and $q\in\mu_V'$, where $\mu_V'$ is
  the equality-mate set of some selected containment witness inside
  $V$.  Symmetric to (ii.b$'$): the selected witness $v'$ satisfies
  $v\eqc v'$ and $u\Reach_R v'$.  Case~(b) with $u'=u$.

\item Subcase (ii.d$'$): $p\in\mu_V$ and $q\in\mu_V'$ as above.  Both
  substitutions apply; the pair $(u',v')$ falls under~(a) by
  composition of (ii.a$'$)--(ii.c$'$).  Case~(b) in full.

\item Subcase (ii.e$'$): both $p$ and $q$ are selected containment
  witnesses of $V$ (neither is $s_V$ or one of its mates).  Then by
  (M1)--(M5) on $V$ and the round-4 argument applied \emph{inside
  $V$}, the source $s_V$ is the common third position with
  $L_M(s_V,p),L_M(s_V,q)\in\{\PP,\PPI\}$ and $L_M(p,q)=R$ forced by
  composition through $s_V$.  Each option places the realising pair
  $(u,v)$ on the vertical backbone of $V$ through the source
  occurrence of $V$, giving case~(a).  Equality-mate variants are
  handled exactly as in (ii.b$'$)--(ii.d$'$).
\end{itemize}

The only difference from the round-4 case analysis is that the
``common stratum'' is no longer forced to be the root stratum $V_0$.
By (M6$'$.b) it is \emph{some} stratum $V\in\Stratum_M$, possibly
nested; the proof goes through inside $V$ because each individual
stratum is itself a round-4-style local mosaic over $P_V$.  In every
subcase, $(u,v)$ satisfies~(a) directly or admits $\eqc$-mate
substitutes that do, which is~(b).
\end{proof}

\begin{remark}
The key insight is that the round-4 Subcase (ii.e) argument is
purely \emph{local} to one stratum: it uses only the round-4 axioms
(M1)--(M5) restricted to $P_V$, plus (M6) restricted to $V$'s
containment data.  Recursive verticalization preserves this locality
inside each stratum and uses (M6$'$) only to ensure that the right
stratum exists.  Cross-stratum pairs never receive $\PP/\PPI$ labels
(by (M6$'$.c)), so the round-4 case analysis is not lost; it is
applied \emph{inside} each stratum.
\end{remark}

%% ==================================================================
\section{Updated validity clauses (V3), (V4), (V5), (V9.Eb)}
\label{sec:validity}
%% ==================================================================

The round-4 validity clauses on certificates (Section~6.5 of
\cite{GPT55Repaired}) are formulated in terms of port-indexed mate
clusters $\Mate_\Q(C,s,p)\subseteq S\times\Port(s')$ attached to
contexts $C\in\Ctx_\Q$.  In the recursive setting the clauses traverse
the stratum forest of each mosaic, in addition to the mate clusters.

We restate only the four clauses affected by the repair.  All other
clauses (V1), (V2), (V6), (V7), (V8), (V9.Ea), (V10) of
\cite{GPT55Repaired} are preserved verbatim.

\subsection*{(V3$'$) Selected containment witness, per stratum.}
For every $\exists R.D\in\Cl(C_0)$ with $R\in\{\PP,\PPI\}$, every
mosaic $M\in\Mos_\Q$, every stratum $V\in\Stratum_M$, and every port
$p$ of the source $s_V$ such that $\exists R.D\in\tau(s_V)$ at port
$p$: there exists a selected $R$-witness $w\in P_V\setminus\mu_V$
declared in $V$, with $D\in\tau(w)$ and $L_M(s_V,w)=R$.
\smallskip\noindent
\emph{Difference from round-4 (V3).} Round-4 (V3) ranged over the
single stratum $\Vrt_M$.  (V3$'$) ranges over every stratum $V$
in the forest $\mathcal V_M$.  Existential discharge for a position
$s_V$ is checked within the stratum that $s_V$ heads.

\subsection*{(V4$'$) Mate-cluster closure, per stratum.}
For every mosaic $M$, every stratum $V\in\Stratum_M$, and every port
$p$ of $s_V$: the port-indexed mate cluster $\Mate_\Q(C_V,s_V,p)$
declared for $V$ is closed under explicit $\EQ$-port chains in $E_\Q$
\emph{within} $V$, in the sense that every chain of $\EQ$-port links
starting from $(s_V,p)$ and remaining inside $P_V$ produces only
ports that lie in $\Mate_\Q(C_V,s_V,p)$.
\smallskip\noindent
\emph{Difference from round-4 (V4).} The closure operation is now
restricted to stratum-local $\EQ$-port chains; cross-stratum
$\EQ$-port chains are forbidden by (M6$'$.a), so they need not be
considered.

\subsection*{(V5$'$) Universal propagation, per stratum.}
For every $\forall R.D\in\Cl(C_0)$, every mosaic $M$, every stratum
$V\in\Stratum_M$, and every pair $p,q\in P_V$ with $L_M(p,q)=R$ and
$\forall R.D\in\tau(p)$: $D\in\tau(q)$.  Furthermore, for
cross-stratum pairs $(p,q)$ with $L_M(p,q)=R\in\{\DR,\PO\}$ and
$\forall R.D\in\tau(p)$: $D\in\tau(q)$ as well (this is the
cross-stratum case, unchanged from round-4).
\smallskip\noindent
\emph{Difference from round-4 (V5).} Within-stratum universal
propagation is restricted to one stratum; cross-stratum $\DR/\PO$
universal propagation is unchanged.  No within-stratum $\PP/\PPI$
universal escapes from $V$ because (M6$'$.b) forbids any $\PP/\PPI$
label between $P_V$ and the rest of the mosaic.

\subsection*{(V9.Eb$'$) Equality-mate joint witnessing, per stratum.}
For every mosaic $M$, every stratum $V\in\Stratum_M$, every explicit
$\EQ$-class $E\subseteq P_V$, and every $\exists R.D\in\Cl(C_0)$
common to the type $\tau(p)$ for all $p\in E$: there exists a single
selected witness $w\in P_V$ with $L_M(p,w)=R$ for some $p\in E$ and
$D\in\tau(w)$, such that the substitution lemma applied to any
$p'\in E$ yields the same $\eqc$-class for $w$.
\smallskip\noindent
\emph{Difference from round-4 (V9.Eb).} Joint witnessing is checked
per stratum, because (M6$'$.a) prevents $\EQ$-classes from spanning
multiple strata.  The clause matches the local structure: an
equality class lives inside one stratum, so its joint discharge is
also stratum-local.

\medskip

\begin{remark}\label{rem:carry-over}
The four clauses above are essentially \emph{universally quantified
over strata}.  Their finite-checkability is preserved: the stratum
forest of every mosaic in $\Mos_\Q$ is finite (a finite set of
strata), so each of (V3$'$)--(V9.Eb$'$) is a finite conjunction of
round-4-style clauses, one per stratum.  The original round-4
validity clauses (V3), (V4), (V5), (V9.Eb) are recovered as the
single-stratum case (depth-1 stratum forest).
\end{remark}

%% ==================================================================
\section{Side-interface width, recurrence over modal depth}
\label{sec:side-bound}
%% ==================================================================

The round-4 paper proves an explicit cubic bound
$|\Side(u)|\leq 1+n^3+3n^2+3n$ for the side interface at an
occurrence $u$, where $n=|C_0|$ (Lemma~4.4 of
\cite{GPT55Repaired}).  Under recursive verticalization, a single
side-witness $w$ of $u$ can itself open a local stratum carrying up
to $|\Side(w)|$ further positions.  The width bound therefore becomes
a recurrence over modal depth.

\subsection{Recurrence}

Let $\md(C_0)\in\mathbb N$ denote the modal depth of $C_0$.  Define
$m:\{0,1,\dots,\md(C_0)\}\to\mathbb N$ by
\[
  m(0)\;=\;1,\qquad
  m(d{+}1)\;=\;1\;+\;n^3\cdot\bigl(1+k_\Mate+k_{\rm bd}\bigr)\;+\;n^3\cdot m(d),
\]
where $n=|C_0|$ and $k_\Mate,k_{\rm bd}$ are the round-4 mate and
boundary constants.

\begin{proposition}[Side-interface width]\label{prop:side-width}
For every occurrence $u$ in $\U(\Q)$ at modal depth $\leq d$ from a
generated context source, $|\Side(u)|\leq m(d)$.
\end{proposition}

\begin{proof}
Induction on $d$.  Base $d=0$: $u$ is a leaf of the witness selector;
$\Side(u)=\{u\}$, so $|\Side(u)|=1=m(0)$.  Inductive step: the
selector at $u$ adds up to $n^3$ selected $\DR/\PO$ witnesses
(round-4 cubic bound), each with up to $1+k_\Mate+k_{\rm bd}$
boundary or equality-mate companions, and each side-witness $w$ may
open a nested stratum $V_w$ of width $\leq m(d)$ (by the induction
hypothesis applied at $w$).  Summing: $|\Side(u)|\leq 1 +
n^3(1+k_\Mate+k_{\rm bd}) + n^3 m(d) = m(d+1)$.
\end{proof}

\begin{corollary}[Effectively computable side-width]\label{cor:m-effective}
$m(\md(C_0))$ is effectively computable from $C_0$.  In particular,
the total side-interface alphabet is finite and bounded by
$m(\md(C_0))\leq (n^3+1)^{\md(C_0)+1}\cdot
\bigl(1+k_\Mate+k_{\rm bd}\bigr)$.
\end{corollary}

The bound $m(\md(C_0))$ is no longer polynomial in $n$ in general; it
is exponential in $\md(C_0)$, which can itself be linear in $n$.  This
matches GPT-5.5's round-5 recommendation: drop the polynomial claim
and commit only to an effectively computable bound.  The certificate
space $B(C_0)$ from Section~6.6 of \cite{GPT55Repaired} now uses
$m(\md(C_0))$ in place of the cubic constant, but the overall
double-exponential bound $B(C_0)\leq 2^{2^{p(n)}}$ for some
computable polynomial $p$ is preserved.

\begin{remark}\label{rem:bound-not-load-bearing}
The decision procedure does not rely on a polynomial side-width.  It
only uses that $B(C_0)$ is effectively computable from $C_0$.  Both
$m(\md(C_0))$ and the round-4 cubic bound satisfy this requirement.
\end{remark}

%% ==================================================================
\section{Worked example: a v2 certificate for $C_{\rm side}$}
\label{sec:worked-example}
%% ==================================================================

We construct a small certificate for
$C_{\rm side}=\exists\DR.(A\sqcap\exists\PP.B)\sqcap\exists\DR.B$
explicitly, to confirm that the v2 discipline accepts the satisfiable
concept.

\subsection{Setup}

Let $C_0=C_{\rm side}$.  The closure is $\Cl(C_0)=\{C_0, \exists\DR.(A
\sqcap\exists\PP.B), A\sqcap\exists\PP.B, A, \exists\PP.B, B,
\exists\DR.B, \top, \bot\}$ together with negations as needed.  Modal
depth $\md(C_0)=2$.

\subsection{Hintikka types}

Three Hintikka types suffice for the source and its two selected
$\DR$-witnesses; a fourth type covers the $\PP$-witness of the first
selected $\DR$-witness:
\begin{align*}
  \tau_x &= \{\,C_0,\,\exists\DR.(A\sqcap\exists\PP.B),
              \exists\DR.B,\,\top\,\},\\
  \tau_y &= \{\,A\sqcap\exists\PP.B,\,A,\,\exists\PP.B,\,\top\,\},\\
  \tau_z &= \{\,B,\,\top\,\},\\
  \tau_w &= \{\,B,\,\top\,\}.
\end{align*}
(We use $\tau_w$ for $y$'s selected $\PP$-witness; it happens to be
equal to $\tau_z$, which is what justifies the identification
$w=z$ in the round-5 stress reading; see Remark~\ref{rem:identify}
below.)

\subsection{Certificate structure}

The certificate $\Q$ has exactly one context $C\in\Ctx_\Q$ whose
source has profile/type $\tau_x$ at port $p_x$.  The mosaic at $C$ is
$M=(P,\tau,L,\mathcal V)$ with $P=\{x,y,z\}$ (we use the same names
for positions and the model elements of
Example~\ref{ex:stress-model}; this is a notational coincidence, not
a semantic identification).  Stratum forest
$\mathcal V=(\Stratum,\Subsrc)$ with:
\begin{itemize}
\item Root stratum $V_0$: $P_{V_0}=\{x\}$, source $s_{V_0}=x$, mates
  $\mu_{V_0}=\{x\}$.  (Source has no explicit equality mates and no
  selected $\PP/\PPI$ witnesses, so $V_0$ contains only $x$.)
\item Nested stratum $V_1$: $P_{V_1}=\{y,z\}$, source $s_{V_1}=z$,
  mates $\mu_{V_1}=\{z\}$, selected $\PPI$-witness $y$.
\item $\Subsrc(V_1)=z$, with $z$ as parent position.  Note
  $z\in P_{V_0}^{\,c}$ in the sense that $z\notin V_0$, but $z$ is a
  sibling-frontier position of $V_0$; (SF5) is satisfied because
  $P_{V_0}\cap P_{V_1}=\emptyset$, $V_1$ being rooted at $z$ which
  itself lies outside $V_0$'s support.
\end{itemize}

Equivalently: the root stratum $V_0$ has source $x$ and no selected
witnesses; the sibling-frontier of $V_0$ contains $y,z$; the nested
stratum $V_1$ rooted at $z$ has $y$ as a $\PPI$-witness of $z$.

\begin{remark}\label{rem:identify}
The original model of Example~\ref{ex:stress-model} has $y\PP z$, so
in the certificate $z$ is the proper-superpart and $y$ the
$\PPI$-mate of $z$; equivalently, $y$ is a selected $\PPI$-witness of
$z$ inside the nested stratum $V_1$.  We do not identify $w$ with $z$
in the certificate; we only use the type-level identification
$\tau_w=\tau_z$ to keep the example compact.
\end{remark}

\subsection{Label matrix and (M6$'$) verification}

The label matrix $L:P^2\to\Base$ on $P=\{x,y,z\}$ is, modulo
reflexive $\EQ$ on the diagonal and inverses on
$L(q,p)=\Inv(L(p,q))$:
\[
  L(x,y)=\DR,\quad L(x,z)=\DR,\quad L(y,z)=\PP.
\]

We verify (M6$'$.a)--(M6$'$.d) on $M$.
\begin{itemize}
\item (M6$'$.a) $\EQ$-closure within strata.  The only $\EQ$-classes
  are the trivial singletons (no explicit equality identifications),
  so the clause holds vacuously.
\item (M6$'$.b) Containment labels are co-vertical.  The only
  $\PP/\PPI$ label is $L(y,z)=\PP$ (and its inverse $L(z,y)=\PPI$).
  We verify that $(y,z)$ is co-vertical: $y\in P_{V_1}$, $z\in
  P_{V_1}$, so $V_1\in\Stratum(y)\cap\Stratum(z)$.  Moreover the
  local containment data of $V_1$ (source $z$, selected
  $\PPI$-witness $y$) matches $L(y,z)=\PP$.
\item (M6$'$.c) Cross-stratum labels are $\DR/\PO$ or $\EQ$.
  $L(x,y)=\DR$ between $x\in P_{V_0}$ and $y\in P_{V_1}$ is
  cross-stratum: $(x,y)$ is not co-vertical (no common stratum
  contains both).  The label $\DR\in\{\DR,\PO\}$ satisfies (M6$'$.c).
  Likewise $L(x,z)=\DR$ between $x\in P_{V_0}$ and $z\in P_{V_1}$:
  cross-stratum, label $\DR$, satisfies (M6$'$.c).
\item (M6$'$.d) Residual sibling frontier.  $x$ is the source of
  $V_0$ and has no other positions in its residual sibling frontier
  (apart from $y,z$ which are handled by (M6$'$.c)).  No other
  residual-frontier pairs exist in $M$.
\end{itemize}
All four clauses hold, so $M$ is a v2-mosaic, and $C_{\rm side}$ is
accepted by the v2 certificate.

\subsection{Validity clauses on this certificate}

We verify the four affected clauses (V3$'$), (V4$'$), (V5$'$),
(V9.Eb$'$).
\begin{itemize}
\item (V3$'$) Selected containment witness.  The only
  $\exists R.D\in\Cl(C_0)$ with $R\in\{\PP,\PPI\}$ is
  $\exists\PP.B\in\tau_y=\tau(y)$.  We check stratum $V_1$: $y$ is a
  position of $V_1$, and $V_1$ declares $z$ as the source with $y$ a
  selected $\PPI$-witness.  $L(y,z)=\PP$ and $B\in\tau(z)$, so $y$'s
  $\exists\PP.B$ is discharged in stratum $V_1$ as required.
\item (V4$'$) Mate-cluster closure.  No explicit $\EQ$-port chains
  are declared, so the clause holds vacuously.
\item (V5$'$) Universal propagation.  No $\forall R.D$ universals are
  in any type of this certificate, so the clause holds vacuously.
\item (V9.Eb$'$) Equality-mate joint witnessing.  No multi-element
  $\EQ$-classes are declared, so the clause holds vacuously.
\end{itemize}
All non-vacuous validity conditions hold.  The unaffected clauses
(V1), (V2), (V6), (V7), (V8), (V9.Ea), (V10) are checked as in
round-4; the only nontrivial point is (V2) (composition consistency),
which holds because $L(x,y)=\DR$, $L(y,z)=\PP$, $L(x,z)=\DR$ and
$\DR\in\comp(\DR,\PP)=\{\DR,\PO,\PP\}$.

\subsection{Confirmation against the repository's reasoners}

The verification script
\texttt{verification/python/wp7\_side\_witness\_stress.py} reports
$C_{\rm side}$ as SAT under all three reasoners.  The v2 certificate
above is the smallest witness-generated certificate matching the
abstract RCC5 model of Example~\ref{ex:stress-model}.

%% ==================================================================
\section{What carries over from round-4 verbatim}
\label{sec:carryover}
%% ==================================================================

The following sections of \cite{GPT55Repaired} are unaffected by the
repair, in the sense that their statements and proofs go through
verbatim once Definition~\ref{def:mosaic-v2}, Lemma~\ref{lem:vertreal-recursive},
the restated validity clauses of \S\ref{sec:validity}, and the
recurrence bound of \S\ref{sec:side-bound} are inlined.

\begin{itemize}
\item Section 2 (Abstract RCC5 frames).
\item Section 3 (Closure, types, and witness thinning), including the
  frame inheritance lemma and the witness-generated submodel theorem.
\item Section 4.1--4.2 (Occurrences, blocking, and equality;
  generated split-cover construction).  Section 4.3 (DR/PO witness
  placement) is amended only at the width bound (Lemma~4.4 of
  \cite{GPT55Repaired}), replaced by Proposition~\ref{prop:side-width}.
\item Section 5.2 (Abstract triple graph $\T_\Q$), with stratum
  membership $\zeta\in\{\mathsf{vertical}_V : V\in\Stratum\}\cup\{\mathsf{sibling}\}$
  ranging over the per-position stratum index instead of a Boolean.
  The bound $|\Pos_\Q|\leq |S|\cdot\max_s|\Port(s)|\cdot|\Ctx_\Q|\cdot
  k_{\rm strat}$, where $k_{\rm strat}$ is the maximum stratum count
  per context (bounded by $m(\md(C_0))$).
\item Section 5.3--5.6 (Renz--Nebel patchwork).  The patchwork lemma
  is applied to the local network on each stratum independently and
  to the residual cross-stratum frontier as in round-4;\@ the
  $\DR/\PO$ arc-consistency argument that the round-4 paper uses on
  the sibling-frontier carries over to the cross-stratum case.
\item Section 6.3 (Port-indexed mate clusters $\Mate_\Q(C,s,p)$),
  unchanged in form;\@ the clusters now attach to per-stratum sources
  and ports, but the definition and (V9.Ea) clause are unchanged.
\item Section 6.5 (Request-fulfilled cycles), unchanged.
\item Section 7.3 (Standalone typed-$\EQ$ quotient theorem),
  unchanged.  The six assumptions (A1)--(A6) remain the same;\@
  (A4) is now discharged by the restated (V3$'$)/(V4$'$)/(V9.Eb$'$)
  combined.
\item Section 8 (Completeness, including the no-automata
  regularization), unchanged.
\end{itemize}

%% ==================================================================
\section{Discussion of the other round-5 items}
\label{sec:other-round5}
%% ==================================================================

GPT-5.5's round-5 review~\cite{GPT55Round5Review} lists eight
recommended repairs.  This paper closes \textbf{Item 2} (side-interface
verticalization discipline), \textbf{Item 3} (add the stress test --
done by Example~\ref{ex:stress-model} and the worked certificate of
\S\ref{sec:worked-example}, with computational corroboration via
\texttt{wp7\_side\_witness\_stress.py}), and \textbf{Item 9} (drop the
polynomial side-width claim, via the recurrence of
\S\ref{sec:side-bound}).  Item~1 (TeX/PDF mismatch) is editorial and
the repository copy is already self-consistent.

The remaining items require independent editorial passes on
\cite{GPT55Repaired}:

\begin{itemize}
\item \textbf{Item~4 (non-circular mosaic patchwork)}: Sections 5.4--5.6
  of \cite{GPT55Repaired} should be reordered so that the abstract
  pair-label function and triple set are defined before the
  patchwork lemma is invoked, and a finite composition-table check on
  every abstract triple is required as a finite validity clause.
  No change to the v2 (M6$'$) statements is needed.
\item \textbf{Item~5 (overlap amalgamation)}: the overlap-amalgamation
  lemma should require a single cross-label assignment compatible
  with every overlap node and every mixed triple, not merely
  per-triple existence.  This is a tightening of the lemma's
  hypothesis, orthogonal to (M6$'$).
\item \textbf{Item~6 (joint equality-witness compatibility as an
  explicit clause)}: the assumption (A4) of the typed-$\EQ$ quotient
  theorem (Section 7 of \cite{GPT55Repaired}) should be promoted to
  an explicit finite validity clause on the certificate.  (V9.Eb$'$)
  of \S\ref{sec:validity} above is the natural per-stratum form;\@
  cross-stratum joint-witnessing is vacuous by (M6$'$.a) (no
  cross-stratum equality classes).
\item \textbf{Item~7 (half-open cycle convention)}: the request-cycle
  lemma in Section 6.5 of \cite{GPT55Repaired} should be rewritten
  with a half-open $[i,j)$ convention and an explicit endpoint
  mapping.  This is orthogonal to (M6$'$).
\item \textbf{Item~8 (global finite regularization)}: the
  no-automata regularization argument in Section 8 of
  \cite{GPT55Repaired} should be promoted to an explicit global
  lemma covering all infinite vertical rays uniformly.  This is
  orthogonal to (M6$'$).
\end{itemize}

Items 4--8 do not affect the (M6$'$) repair and can be addressed in
independent editorial passes.  The v2 architecture of this paper is
the natural setting for those passes, since each of them is local to
a specific section of \cite{GPT55Repaired}.

%% ==================================================================
\section{Summary}
%% ==================================================================

The repair is a recursive lift of the round-4 verticalization
discipline.  The only structural change is that the flat
$\Vrt\subseteq P$ of a mosaic is replaced by a finite \emph{stratum
forest} $\mathcal V$, in which selected $\DR/\PO$ side-witnesses
may themselves open nested local vertical strata.  The repair:
\begin{enumerate}[label=(\arabic*)]
\item closes the round-5 stress concept
  $\exists\DR.(A\sqcap\exists\PP.B)\sqcap\exists\DR.B$ by accepting
  the small certificate of \S\ref{sec:worked-example};
\item preserves Lemma~7.1 of \cite{GPT55Repaired} (vertical
  realisation of $\PP/\PPI$ pairs) by the recursive case analysis of
  \S\ref{sec:lifting};
\item preserves the round-4 case analysis of Subcase (ii.e) of the
  Vertical-Realisation Lemma by restricting the analysis to one
  stratum at a time;
\item recovers the round-4 (M6) discipline as the depth-1 case of
  (M6$'$);
\item preserves the rest of the round-4 architecture verbatim,
  modulo the per-stratum quantifier on (V3), (V4), (V5), (V9.Eb);
\item replaces the polynomial side-width $m(n)=O(n^3)$ with an
  effectively computable recurrence over modal depth, which suffices
  for decidability;\@ the double-exponential certificate-space bound
  $B(C_0)\leq 2^{2^{p(n)}}$ is preserved.
\end{enumerate}

The repair is computationally corroborated: all three Python
reasoners in the repository (cover-tree tableau, baseline and
cycle-aware quasimodel) report $C_{\rm side}$ as SAT.  The v2
discipline (M6$'$) accepts that satisfiable concept by construction.

%% ==================================================================
\begin{thebibliography}{99}

\bibitem{GPT55Repaired}
GPT-5.5 (OpenAI), prompted by Michael Wessel.
\emph{A Generated Split-Forest Decidability Proof for $\ALCIRCC{5}$
Without Automata: A Repaired Containment-Oriented Presentation}.
Unpublished manuscript, May 2026.  Repository:
\url{https://github.com/lambdamikel/alcircc5/blob/master/papers/gpt5.5_round2/repaired_split_forest_no_automata_proof.pdf}.

\bibitem{GPT55Round5Review}
GPT-5.5 (OpenAI), prompted by Michael Wessel.
\emph{Formal Review of the Latest Repaired Split-Forest Proof
Manuscript}.  Unpublished manuscript, May~22, 2026.  Repository:
\url{https://github.com/lambdamikel/alcircc5/blob/master/papers/gpt5.5_round2/formal_review_progress_paper.pdf}.

\bibitem{SiblingV1}
GPT-5.4~Pro (OpenAI), prompted by Michael Wessel.
\emph{Sibling-Interface Descriptors for $\ALCIRCC{5}$}.  Unpublished
manuscript, April~2026.  Repository:
\url{https://github.com/lambdamikel/alcircc5/blob/master/papers/trees/sibling_interface_descriptors_ALCIRCC5_completed_eqsync_canonical_needpatched.pdf}.

\bibitem{Wessel2003}
Michael Wessel.
\emph{Some Practical Issues in Building Description Logic Based
Spatially-Aware Information Systems}.  In:\@ Description Logics, 2003.

\bibitem{RenzNebel1999}
Jochen Renz and Bernhard Nebel.
\emph{On the Complexity of Qualitative Spatial Reasoning: A Maximal
Tractable Fragment of the Region Connection Calculus}.  Artificial
Intelligence, 108(1--2):69--123, 1999.

\bibitem{ALCIRCC5Repo}
Michael Wessel et al.
\emph{ALCI\_RCC5 Project Repository}.
\url{https://github.com/lambdamikel/alcircc5}.

\end{thebibliography}

\end{document}
